\section{Proofs}\label{app:proofs}

This appendix proves the propositions of \cref{sec:model}. Notation is as in the text: market clearing gives
$m=(u-a)/\Phi$ \eqref{eq:clear}, and the arbitrage sector maximises \eqref{eq:objective}.

\paragraph{Proposition 1.} The objective \eqref{eq:objective} is strictly concave in $a$, with a kink at zero. Its best
response to a basis $m$ is $a=(m-c_+)/\gamma$ if $m>c_+$, $a=(m+c_-)/\gamma$ if $m<-c_-$, and $a=0$ otherwise. In
equilibrium $m=(u-a)/\Phi$. The map $a\mapsto \mathrm{BR}\big((u-a)/\Phi\big)-a$ is strictly decreasing, so the fixed point
is unique. Substituting the clearing condition into each branch of the best response gives \eqref{eq:band}. For the
upper branch, $m\Phi=u-(m-c_+)/\gamma$, so $m=(\gamma u+c_+)/(\gamma\Phi+1)$, which equals $c_+$ at $u=\Phi c_+$, so the
basis is continuous at the edge. The lower branch is symmetric. Inside the band $a=0$ and $m=u/\Phi$.

\paragraph{Proposition 2.} (i) and (ii) follow from \eqref{eq:band}: the edges, $c_+$ and $-c_-$, do not contain
$\gamma$, and the slope outside the band, $\gamma/(\gamma\Phi+1)$, is increasing in $\gamma$. (iii) The costs can be
written as $c_+=\tau_++(r_A-r)(T-t)$ and $c_-=\tau_-+(\ell_s+r-r_L)(T-t)$, where $\tau_\pm$ are the transaction costs of
each trade, $r_A$ is the return required on arbitrage capital, $\ell_s$ is the fee for borrowing the stock and $r_L$ the
interest earned on the proceeds of a short sale. Both edges are therefore linear in the time to expiry $T-t$ and converge
to the transaction costs at expiry. (iv) The midpoint of $[-c_-,c_+]$ is $(c_+-c_-)/2$, which is negative when $c_->c_+$.

\paragraph{Proposition 3.} With an existing position $a_{t-1}\ge 0$, the sector maximises
$a_tm_t-\tfrac{\gamma}{2}a_t^2-c_+(a_t-a_{t-1})^+-c_u(a_{t-1}-a_t)^+$. This is strictly concave with a kink at
$a_{t-1}$. The marginal value of holding one more unit is $m_t-\gamma a_t$, so adding pays only if
$m_t-\gamma a_{t-1}>c_+$ and unwinding only if $m_t-\gamma a_{t-1}<-c_u$; otherwise the sector holds $a_{t-1}$. When it
adds, the interior solution is $a_t=(m_t-c_+)/\gamma$; when it unwinds, $a_t=\max\{0,(m_t+c_u)/\gamma\}$. Combining each
case with the clearing condition gives a unique equilibrium, as in Proposition 1.

\paragraph{Proposition 4.} In a ban, restriction (B1) removes the option to add. Starting from $a_{t-1}>0$, the sector
holds its position as long as $m_t\ge\gamma a_{t-1}-c_u$. In this region $m_t=(u_t-a_{t-1})/\Phi$, so every increase
in $u_t$ passes through to the basis at the rate $1/\Phi$, instead of the smaller rate $\gamma/(\gamma\Phi+1)$ that
applies above the upper edge in normal times: rich futures are not corrected. If $u_t$ is low enough that the sector
unwinds, then $a_t=(m_t+c_u)/\gamma$ when interior, so $m_t=\gamma a_t-c_u$: cheap futures are corrected, at a floor that
falls as the position shrinks. When $a_t$ reaches zero, the floor disappears and $m_t=u_t/\Phi$. Positions that reach
expiry during a ban are delivered and cannot be opened again in the next contract, so they leave the book for the rest of
the ban.

\paragraph{Proposition 5.} Split outside traders' demand shocks into closing demand $\varepsilon^c$ and demand for new
futures exposure $\varepsilon^n\ge 0$. In normal times both enter $u$ one for one. In a ban, by (B2), new demand
cannot be executed in futures; a share $\chi$ of it buys the stock instead, which in the units of \eqref{eq:clear} enters
$u$ as $-\lambda\varepsilon^n$ with $\lambda=\chi\psi_F/\psi_S\ge0$. Compare the basis with and without a ban for the same
starting position, the same initial pressure $u^0$ and the same shocks: $m^N=m(u^0+\varepsilon^c+\varepsilon^n)$ and
$m^B=m^{\text{ban}}(u^0+\varepsilon^c-\lambda\varepsilon^n)$. Since $m^N$ is non-decreasing and $m^B$ non-increasing in
$\varepsilon^n$, the difference $m^B-m^N$ is non-increasing in $\varepsilon^n$. At $\varepsilon^n=0$ the two inputs are
equal, and since (B1) only removes the option to add, which lowers the basis when used, $m^B\ge m^N$. As $\varepsilon^n$
grows, $m^N$ rises without bound (its slope is at least $\gamma/(\gamma\Phi+1)>0$) while $m^B$ does not rise, so
$m^B<m^N$ for large enough $\varepsilon^n$. The sign of the ban's effect on the level therefore depends on how much of the
pressure is new demand rather than closing demand.

\paragraph{Proposition 6.} If every trader may only reduce positions in each contract, a buyer can only be closing a
short and a seller only closing a long, so every trade lowers open interest by its size: $\Delta OI=-V$ outside expiry
days, and $R=(\Delta OI+V)/V=0$. Trades require that some holders of long positions and some holders of short positions
both want to close at the same price, so volume falls. If an entity may instead lower its overall delta by opening an
offsetting contract in another instrument, as allowed from October 2025, a trade can open a position in one contract,
and the identity no longer holds.
